% ============================================================================
%  Memoria del Trabajo Fin de Grado — ComprendiA
%  Facultad de Informática · Universidad Pontificia de Salamanca
%
%  Formato según NormasEstructuraEstiloMemoriaTFG:
%   - A4, doble cara
%   - Márgenes: superior/inferior 3 cm, interno 3,5 cm, externo 2,5 cm
%   - Preliminares en numeración romana; memoria en arábiga desde 1
%   - Encabezado: impares = capítulo; pares = título del TFG
%
%  Compilar:  pdflatex main  →  bibtex main  →  pdflatex main x2
% ============================================================================
\documentclass[12pt,a4paper,twoside,openright]{report}

\usepackage[utf8]{inputenc}
\usepackage[T1]{fontenc}
\usepackage[spanish,es-tabla]{babel}
\usepackage{lmodern}
\usepackage{microtype}

% --- Márgenes exigidos por la normativa -------------------------------------
\usepackage[a4paper,top=3cm,bottom=3cm,inner=3.5cm,outer=2.5cm]{geometry}

\usepackage{graphicx}
\usepackage{float}
\usepackage{booktabs}
\usepackage{longtable}
\usepackage{enumitem}
\usepackage{xcolor}
\usepackage{listings}
\usepackage{csquotes}
\usepackage[hidelinks]{hyperref}
\usepackage{fancyhdr}
\usepackage{titlesec}
\usepackage{setspace}

\onehalfspacing

% --- Estilo de código -------------------------------------------------------
\definecolor{codegray}{gray}{0.95}
\lstset{
  basicstyle=\ttfamily\footnotesize,
  backgroundcolor=\color{codegray},
  breaklines=true,
  frame=single,
  showstringspaces=false,
  numbers=left,
  numberstyle=\tiny,
  captionpos=b
}

% --- Encabezados: impar = capítulo, par = título del TFG --------------------
\newcommand{\tituloTFG}{ComprendiA: copiloto educativo para vídeos de YouTube}
\pagestyle{fancy}
\fancyhf{}
\fancyhead[LE]{\small\tituloTFG}              % pares: título del TFG
\fancyhead[RO]{\small\nouppercase{\leftmark}} % impares: capítulo
\fancyfoot[LE,RO]{\thepage}                    % paginación exterior
\renewcommand{\headrulewidth}{0.4pt}
\fancypagestyle{plain}{\fancyhf{}\fancyfoot[LE,RO]{\thepage}\renewcommand{\headrulewidth}{0pt}}

\begin{document}

% ====================  PRELIMINARES (numeración romana)  ====================
\frontmatter
\pagenumbering{roman}

% ---------------------------------------------------------------------------
%  Portada y primera página (con firma del director).
%  La normativa pide: escudo UPSA, "Trabajo Fin de Grado" + título completo,
%  estudiante, director, "Salamanca, mes de año".
% ---------------------------------------------------------------------------
\begin{titlepage}
  \centering
  % TODO: sustituir por el escudo oficial de la UPSA (escudo-upsa.png en /img)
  % \includegraphics[width=0.28\textwidth]{img/escudo-upsa.png}\\[1.2cm]
  {\large Universidad Pontificia de Salamanca\\Facultad de Informática\par}
  \vfill
  {\Large\bfseries Trabajo Fin de Grado\par}
  \vspace{1.5cm}
  {\huge\bfseries ComprendiA:\\[0.4cm] copiloto educativo para el análisis y consulta de vídeos de YouTube\par}
  \vspace{2cm}
  {\large Grado en Ingeniería Informática\par}
  \vfill
  \begin{flushright}
    \large
    \textbf{Autor:} Iván Zheng\\[0.4cm]
    \textbf{Director:} % TODO: nombre del director/a
  \end{flushright}
  \vfill
  {\large Salamanca, junio de 2026\par} % TODO: ajustar mes/año de entrega
\end{titlepage}
\thispagestyle{empty}

% --- Primera página interior: igual que la portada, con firma del director ---
\cleardoublepage
\begin{titlepage}
  \centering
  {\large Universidad Pontificia de Salamanca\\Facultad de Informática\par}
  \vfill
  {\Large\bfseries Trabajo Fin de Grado\par}
  \vspace{1.5cm}
  {\huge\bfseries ComprendiA:\\[0.4cm] copiloto educativo para el análisis y consulta de vídeos de YouTube\par}
  \vspace{2cm}
  {\large Grado en Ingeniería Informática\par}
  \vfill
  \begin{flushright}
    \large
    \textbf{Autor:} Iván Zheng\\[1.2cm]
    \rule{6cm}{0.4pt}\\[0.2cm]      % línea para la firma
    \textbf{Vº Bº del Director:} % TODO: nombre del director/a
  \end{flushright}
  \vfill
  {\large Salamanca, junio de 2026\par}
\end{titlepage}
\thispagestyle{empty}

% --- Página en blanco ---
\cleardoublepage
\thispagestyle{empty}
\mbox{}


% ---------------------------------------------------------------------------
%  Resumen (200-250 palabras), Abstract y Descriptores.
% ---------------------------------------------------------------------------
\cleardoublepage
\chapter*{Resumen}
\addcontentsline{toc}{chapter}{Resumen}

ComprendiA es una aplicación web educativa que transforma vídeos de YouTube en material
de estudio consultable. A partir de la URL de una grabación, el sistema obtiene su
transcripción —reutilizando los subtítulos de YouTube cuando existen y recurriendo al
modelo Whisper en caso contrario—, la divide en fragmentos y genera un \textit{embedding}
por fragmento que permite la búsqueda semántica sobre el contenido. Sobre esa base, un
proceso de análisis automático produce capítulos navegables y conceptos clave con sus
marcas de tiempo, y un asistente conversacional responde preguntas del alumno mediante una
arquitectura de generación aumentada por recuperación (RAG), manteniendo una memoria corta
del diálogo para resolver referencias implícitas.

La aplicación organiza las clases procesadas en asignaturas. Para reducir el trabajo manual,
incorpora una clasificación automática que sugiere a qué asignatura pertenece cada vídeo
nuevo, primero por coincidencia de canal de YouTube y, si no la hay, por similitud semántica
con las asignaturas existentes, creando una nueva cuando ninguna resulta adecuada.

El sistema se ha construido con un \textit{backend} en Java 21 sobre Quarkus, una base de
datos PostgreSQL con la extensión \texttt{pgvector} y un \textit{frontend} en Angular. Esta
memoria describe la motivación, el estado del arte de las tecnologías implicadas, la
especificación de requisitos, el diseño y la implementación del sistema, las pruebas
realizadas y las líneas de trabajo futuras.

\vspace{1cm}
\section*{Abstract}

% TODO: revisar la traducción con un hablante nativo antes de la entrega.
ComprendiA is an educational web application that turns YouTube videos into searchable study
material. Given a video URL, the system obtains its transcript —reusing YouTube subtitles when
available and falling back to the Whisper model otherwise—, splits it into fragments and
generates one embedding per fragment to enable semantic search over the content. On top of
this, an automatic analysis produces navigable chapters and key concepts with their
timestamps, and a conversational assistant answers the student's questions using a
retrieval-augmented generation (RAG) architecture, keeping a short-term memory of the dialogue
to resolve implicit references.

The application organises processed lectures into subjects. To reduce manual effort, it
includes an automatic classifier that suggests which subject each new video belongs to, first
by matching the YouTube channel and, failing that, by semantic similarity with existing
subjects, creating a new one when none fits.

The system is built with a Java 21 backend on Quarkus, a PostgreSQL database with the
\texttt{pgvector} extension and an Angular frontend. This document covers the motivation,
the state of the art, the requirements specification, the design and implementation, the
testing performed and future work.

\vspace{1cm}
\section*{Descriptores / Keywords}

\noindent\textbf{Descriptores:} transcripción automática, búsqueda semántica, \textit{embeddings},
generación aumentada por recuperación (RAG), tecnología educativa.

\noindent\textbf{Keywords:} automatic transcription, semantic search, embeddings,
retrieval-augmented generation (RAG), educational technology.
   % resumen, abstract, descriptores

\cleardoublepage
\tableofcontents
\cleardoublepage
\listoffigures
\cleardoublepage
\listoftables

% ====================  MEMORIA (numeración arábiga)  ========================
\mainmatter
\pagenumbering{arabic}

\chapter{Introducción}
\label{cap:introduccion}

\section{Justificación y motivación}

El vídeo se ha consolidado como uno de los principales soportes de aprendizaje. Plataformas
como YouTube concentran una enorme cantidad de clases, charlas y tutoriales, pero el formato
audiovisual presenta una limitación clara para el estudio: es lineal y poco consultable. Para
localizar un concepto concreto dentro de una grabación de una hora, el estudiante debe recordar
aproximadamente en qué momento se trató y desplazarse manualmente por la barra de
reproducción. A diferencia de un texto, un vídeo no se puede buscar, ni resumir, ni hojear con
facilidad.

Existen herramientas parciales —generación de subtítulos, capítulos manuales o resúmenes
automáticos—, pero rara vez se integran en un único flujo pensado para el estudio. Lo habitual
es que el alumno combine varias aplicaciones y siga sin poder \textit{preguntar} directamente
al contenido del vídeo.

ComprendiA nace para cubrir ese hueco: convertir cualquier vídeo de YouTube en material de
estudio estructurado y consultable. La aplicación transcribe la grabación, la analiza para
extraer capítulos y conceptos clave, y ofrece un asistente conversacional capaz de responder
preguntas sobre el contenido citando el momento exacto del vídeo. Además, organiza las clases
en asignaturas y reduce el trabajo administrativo sugiriendo automáticamente a qué asignatura
pertenece cada nueva grabación.

\section{Objetivos}

\subsection{Objetivo general}

Diseñar e implementar una aplicación web que permita transformar vídeos de YouTube en material
de estudio estructurado, ofreciendo transcripción, análisis automático del contenido y consulta
mediante un asistente conversacional basado en recuperación semántica.

\subsection{Objetivos específicos}

\begin{itemize}
  \item Obtener la transcripción de un vídeo de YouTube de forma eficiente, priorizando los
        subtítulos existentes y recurriendo a un modelo de reconocimiento del habla cuando sea
        necesario.
  \item Indexar el contenido mediante \textit{embeddings} para permitir búsqueda semántica por
        fragmentos.
  \item Generar automáticamente capítulos navegables y conceptos clave con sus marcas de tiempo.
  \item Implementar un asistente conversacional (RAG) que responda preguntas sobre el vídeo con
        memoria corta del diálogo y referencias al momento concreto de la grabación.
  \item Organizar las clases en asignaturas y sugerir automáticamente la asignatura de cada
        vídeo nuevo mediante criterios de canal y de similitud semántica.
  \item Construir una interfaz web clara que integre reproductor, capítulos, conceptos y chat.
\end{itemize}

\section{Alcance del proyecto}

El proyecto abarca el desarrollo completo de la aplicación, tanto del \textit{backend} (API,
procesamiento y persistencia) como del \textit{frontend} (interfaz de usuario). Se contempla el
procesamiento de vídeos públicos de YouTube y la gestión de clases y asignaturas por parte de un
único perfil de usuario.

Quedan fuera del alcance de esta primera versión la gestión multiusuario con autenticación, el
procesamiento de fuentes de vídeo distintas a YouTube y el despliegue en un entorno de
producción con alta disponibilidad. Estos puntos se recogen como líneas futuras en el
capítulo~\ref{cap:conclusiones}.

\section{Estructura de la memoria}

El resto de la memoria se organiza como sigue. El capítulo~\ref{cap:estado-arte} revisa el
estado del arte de las tecnologías implicadas (transcripción automática, \textit{embeddings},
búsqueda semántica y modelos de lenguaje). El capítulo~\ref{cap:requisitos} recoge la
especificación de requisitos del sistema. El capítulo~\ref{cap:desarrollo} describe la
metodología, la arquitectura y la implementación. El capítulo~\ref{cap:pruebas} presenta las
pruebas realizadas y sus resultados. Por último, el capítulo~\ref{cap:conclusiones} expone las
conclusiones y las líneas de trabajo futuras.

\chapter{Estado del arte}
\label{cap:estado-arte}

Este capítulo revisa las tecnologías sobre las que se apoya ComprendiA y las alternativas
consideradas en cada caso, justificando las decisiones tomadas.

\section{Transcripción automática del habla}

La conversión de voz a texto (\textit{Automatic Speech Recognition}, ASR) es el primer paso del
sistema. Se valoraron dos vías:

\begin{itemize}
  \item \textbf{Subtítulos de YouTube.} Muchos vídeos disponen de subtítulos, ya sean manuales o
        generados automáticamente por la plataforma. Reutilizarlos evita descargar y procesar el
        audio, reduciendo notablemente el tiempo y el coste.
  \item \textbf{Whisper} \cite{radford2022whisper}. Modelo de reconocimiento del habla de OpenAI,
        robusto frente a ruido y con buen rendimiento multilingüe. Se emplea como alternativa
        cuando el vídeo no ofrece subtítulos válidos.
\end{itemize}

ComprendiA adopta una estrategia híbrida: intenta primero obtener los subtítulos mediante
\texttt{yt-dlp} y solo recurre a Whisper si no existen o son de calidad insuficiente. Esta
decisión prioriza el rendimiento sin renunciar a la cobertura.

\section{Representación semántica: \textit{embeddings}}

Para poder buscar por significado y no solo por coincidencia de palabras, el texto se
transforma en vectores numéricos (\textit{embeddings}) que sitúan fragmentos de contenido
similar cerca en el espacio vectorial. Se utiliza el modelo \texttt{text-embedding-3-small} de
OpenAI por su equilibrio entre calidad y coste. La similitud entre vectores se mide con la
distancia del coseno.

El almacenamiento y la consulta eficiente de estos vectores se resuelven con la extensión
\texttt{pgvector} \cite{pgvector} de PostgreSQL, que añade un tipo de dato vectorial y
operadores de distancia indexables a una base de datos relacional convencional, evitando
introducir una base de datos vectorial dedicada.

\section{Búsqueda semántica y generación aumentada por recuperación (RAG)}

La arquitectura RAG \cite{lewis2020rag} combina la recuperación de información con la generación
de lenguaje natural: ante una pregunta, primero se recuperan los fragmentos más relevantes
mediante búsqueda semántica y después se entregan a un modelo de lenguaje como contexto para que
redacte la respuesta. Frente al uso directo de un modelo generativo, RAG reduce las alucinaciones
y ancla las respuestas en el contenido real del vídeo, además de permitir citar el momento exacto
de la grabación.

\section{Modelos de lenguaje generativos}

Para la redacción de respuestas, la generación de resúmenes y la extracción de capítulos y
conceptos se emplea \texttt{gpt-4o-mini} de OpenAI, un modelo de coste contenido y latencia
razonable, suficiente para las tareas planteadas. El sistema está diseñado para degradar de
forma controlada: si el modelo no está disponible, se aplican alternativas locales (por ejemplo,
una segmentación por bloques temporales) para no interrumpir el procesamiento.

\section{Herramientas y plataformas similares}

% TODO: ampliar con un breve estudio comparativo (p. ej. funciones de resumen de vídeo de
% terceros, extensiones de navegador, plataformas LMS) y una tabla comparativa.
Existen soluciones que abordan partes del problema —generadores de resúmenes de vídeo,
extensiones que transcriben YouTube o asistentes de estudio genéricos—, pero no integran en un
único flujo la transcripción, el análisis estructurado (capítulos y conceptos), la consulta
conversacional anclada en el vídeo y la organización por asignaturas. ComprendiA se diferencia
precisamente por esa integración orientada al estudio.

\chapter{Especificación de requisitos}
\label{cap:requisitos}

Este capítulo describe qué debe hacer el sistema, con independencia de cómo se implementa.
Se identifican los actores, los requisitos funcionales, los requisitos no funcionales y los
casos de uso principales.

\section{Actores}

\begin{itemize}
  \item \textbf{Estudiante (usuario).} Único actor humano. Procesa vídeos, consulta clases,
        organiza asignaturas e interactúa con el asistente. En esta versión no existe gestión
        multiusuario, por lo que se asume un único perfil.
  \item \textbf{Servicios externos.} YouTube (origen del vídeo y subtítulos, vía \texttt{yt-dlp})
        y la API de OpenAI (Whisper, \textit{embeddings} y modelo de lenguaje). No son actores en
        el sentido clásico, pero condicionan el comportamiento del sistema.
\end{itemize}

\section{Requisitos funcionales}

\begin{longtable}{p{2cm}p{11.5cm}}
  \toprule
  \textbf{ID} & \textbf{Descripción} \\
  \midrule
  \endhead
  RF-01 & El sistema debe aceptar la URL de un vídeo de YouTube y procesarlo de forma asíncrona,
          informando del estado (descarga, transcripción, guardado, \textit{embeddings}, análisis). \\
  RF-02 & El sistema debe obtener la transcripción del vídeo, reutilizando los subtítulos de
          YouTube si existen y recurriendo a Whisper en caso contrario. \\
  RF-03 & El sistema debe dividir la transcripción en fragmentos y generar un \textit{embedding}
          por fragmento. \\
  RF-04 & El sistema debe generar automáticamente capítulos navegables y conceptos clave, cada
          uno con su marca de tiempo. \\
  RF-05 & El usuario debe poder reproducir el vídeo y saltar a un instante concreto desde los
          capítulos, conceptos o respuestas del asistente. \\
  RF-06 & El sistema debe permitir realizar búsquedas semánticas sobre el contenido del vídeo. \\
  RF-07 & El asistente debe responder preguntas sobre el vídeo mediante RAG, manteniendo una
          memoria corta del diálogo para resolver referencias implícitas. \\
  RF-08 & El usuario debe poder editar los metadatos de la clase (asignatura, profesor, fecha) y
          marcarla como completada. \\
  RF-09 & El sistema debe organizar las clases en asignaturas y permitir crear, editar y eliminar
          asignaturas. \\
  RF-10 & Al terminar el análisis, el sistema debe sugerir automáticamente una asignatura para el
          vídeo (por canal de YouTube o por similitud semántica, creando una nueva si no hay
          candidata) y marcarla como sugerencia editable. \\
  RF-11 & El usuario debe poder cambiar manualmente la asignatura sugerida, momento en el que deja
          de considerarse sugerencia. \\
  RF-12 & El sistema debe permitir eliminar una clase y todos sus datos asociados. \\
  \bottomrule
\end{longtable}

\section{Requisitos no funcionales}

\begin{longtable}{p{2cm}p{11.5cm}}
  \toprule
  \textbf{ID} & \textbf{Descripción} \\
  \midrule
  \endhead
  RNF-01 & \textbf{Rendimiento.} El procesamiento debe priorizar los subtítulos sobre Whisper para
           reducir tiempo y coste. La interfaz no debe bloquearse durante el procesamiento. \\
  RNF-02 & \textbf{Robustez.} Si un servicio externo (OpenAI, \textit{yt-dlp}) falla, el sistema
           debe degradar de forma controlada sin interrumpir el procesamiento. \\
  RNF-03 & \textbf{Usabilidad.} La interfaz debe ser clara, en modo oscuro, integrando reproductor,
           capítulos, conceptos y chat en una sola pantalla. \\
  RNF-04 & \textbf{Mantenibilidad.} Todo el código (clases, métodos, variables y comentarios) se
           escribe en español, con una separación clara entre capas. \\
  RNF-05 & \textbf{Coste.} Se emplean modelos de coste contenido (\texttt{gpt-4o-mini},
           \texttt{text-embedding-3-small}) y se evita recalcular \textit{embeddings} ya disponibles. \\
  RNF-06 & \textbf{Portabilidad de datos.} Los \textit{embeddings} se almacenan en PostgreSQL con
           \texttt{pgvector}, sin introducir una base de datos vectorial adicional. \\
  \bottomrule
\end{longtable}

\section{Casos de uso}

% TODO: añadir el diagrama de casos de uso (UML) en img/casos-uso.png
La figura de casos de uso (pendiente) resume las interacciones del estudiante con el sistema. Los
casos de uso principales son: \textit{Procesar vídeo}, \textit{Consultar clase}, \textit{Preguntar
al asistente}, \textit{Buscar en el vídeo}, \textit{Gestionar asignaturas} y \textit{Editar
metadatos de clase}.

\subsection{Caso de uso detallado: Procesar vídeo}

\begin{itemize}
  \item \textbf{Actor:} Estudiante.
  \item \textbf{Precondición:} Disponer de la URL de un vídeo público de YouTube.
  \item \textbf{Flujo principal:}
    \begin{enumerate}
      \item El estudiante introduce la URL y solicita el análisis.
      \item El sistema valida la URL y lanza el procesamiento asíncrono.
      \item El sistema obtiene la transcripción, genera fragmentos y \textit{embeddings}, y produce
            capítulos y conceptos.
      \item El sistema sugiere una asignatura para la clase.
      \item Al completarse, se abre la pantalla de la clase.
    \end{enumerate}
  \item \textbf{Flujos alternativos:} URL no válida (se informa del error); cancelación por el
        usuario (se eliminan los datos parciales); fallo de un servicio externo (se aplica la
        alternativa local correspondiente).
  \item \textbf{Postcondición:} La clase queda persistida y consultable.
\end{itemize}

\chapter{Desarrollo}
\label{cap:desarrollo}

Este capítulo describe cómo se ha construido ComprendiA: la metodología seguida, las
tecnologías empleadas, la arquitectura del sistema, el modelo de datos, el pipeline de
procesamiento y la implementación de cada módulo funcional.

\section{Metodología}

El desarrollo se ha llevado a cabo con una metodología iterativa e incremental, aplicando
principios de Scrum. El trabajo se organizó en \textit{sprints} cortos (de una a dos semanas):
en cada uno se planificaban tareas concretas, se implementaban, se probaban y se revisaban con
el tutor. Este enfoque permitió validar funcionalidades de forma temprana y reorientar el
alcance cuando fue necesario, en lugar de fijar todos los requisitos al inicio.

Cada incremento añadió una capacidad completa y utilizable: primero la transcripción y la
persistencia, después la búsqueda semántica y el asistente, más tarde el análisis automático
(capítulos y conceptos), la organización por asignaturas y, por último, la clasificación
automática de asignaturas.

\section{Tecnologías y herramientas}

\begin{itemize}
  \item \textbf{Backend:} Java~21 sobre \textbf{Quarkus}, con Hibernate ORM (Panache) para la
        persistencia, \textbf{WebSockets} (\texttt{quarkus-websockets-next}) para el progreso en
        tiempo real y ejecución del procesamiento en hilos virtuales de Java~21.
  \item \textbf{Base de datos:} \textbf{PostgreSQL} (alojada en Neon) con la extensión
        \textbf{pgvector} para almacenar y consultar \textit{embeddings}.
  \item \textbf{Frontend:} \textbf{Angular} en modo \textit{zoneless}, con una interfaz en modo
        oscuro que integra reproductor, capítulos, conceptos y chat.
  \item \textbf{Servicios de IA (OpenAI):} \texttt{whisper} (transcripción),
        \texttt{text-embedding-3-small} (\textit{embeddings}) y \texttt{gpt-4o-mini} (chat,
        resúmenes y análisis).
  \item \textbf{Obtención de vídeo:} \texttt{yt-dlp} para descargar audio, subtítulos y
        metadatos del canal, y \texttt{ffmpeg} para el tratamiento del audio.
\end{itemize}

% TODO: justificar brevemente la elección de Quarkus frente a Spring Boot, y de Angular zoneless.

\section{Arquitectura general}

% TODO: incluir el diagrama de arquitectura en img/arquitectura.png
El sistema sigue una arquitectura cliente-servidor. El \textit{frontend} Angular consume una API
REST expuesta por el \textit{backend} Quarkus. El \textit{backend} se organiza en capas:
recursos REST (\texttt{recurso}), servicios de dominio (\texttt{servicio}), repositorios de
acceso a datos (\texttt{repositorio}) y entidades (\texttt{entidad}), además de objetos de
transferencia (\texttt{dto}). Los servicios externos (YouTube vía \texttt{yt-dlp} y la API de
OpenAI) se invocan desde la capa de servicios.

\section{Modelo de datos}

Las entidades principales son \texttt{Video} (una clase procesada), \texttt{Asignatura},
\texttt{Profesor}, \texttt{FragmentoTranscripcion} (con su \textit{embedding}),
\texttt{CapituloVideo} y \texttt{ConceptoClaveVideo}. La tabla~\ref{tab:entidades} resume las
relaciones; el detalle de campos se encuentra en el repositorio del proyecto.

\begin{table}[H]
  \centering
  \caption{Entidades principales del modelo de datos.}
  \label{tab:entidades}
  \begin{tabular}{p{3.2cm}p{9cm}}
    \toprule
    \textbf{Entidad} & \textbf{Descripción y relaciones} \\
    \midrule
    \texttt{Video} & Clase procesada. N--1 con \texttt{Asignatura} y \texttt{Profesor}. Guarda
      metadatos editables y los campos de clasificación automática. \\
    \texttt{Asignatura} & Agrupa vídeos. Incluye datos para la clasificación (canal de YouTube,
      palabras clave y \textit{embedding} representativo). \\
    \texttt{Profesor} & Docente asociado a vídeos y asignaturas. \\
    \texttt{FragmentoTranscripcion} & Fragmento de transcripción con tiempos y \textit{embedding}
      (\texttt{pgvector}). Base de la búsqueda semántica. \\
    \texttt{CapituloVideo} & Capítulo navegable (IA o manual) con tiempos. \\
    \texttt{ConceptoClaveVideo} & Concepto clave con definición y marca de tiempo. \\
    \bottomrule
  \end{tabular}
\end{table}

Para no perder datos antiguos, los nombres de asignatura y profesor se conservan también como
texto y se migran al modelo relacional al arrancar la aplicación.

\section{Pipeline de procesamiento asíncrono}

Al solicitar el análisis de un vídeo, el \textit{backend} responde de inmediato con un
identificador de trabajo y ejecuta el procesamiento en segundo plano. El progreso se transmite
al \textit{frontend} \textbf{en tiempo real mediante WebSockets}: el cliente se conecta al canal
\texttt{/ws/trabajos/\{id\}} y recibe cada cambio de fase según ocurre. Si la conexión WebSocket
no está disponible, el sistema recurre automáticamente a un sondeo (\textit{polling}) HTTP
periódico, de modo que la funcionalidad se mantiene en cualquier caso. La barra de progreso
recorre las fases:

\begin{center}
\texttt{DESCARGANDO} $\rightarrow$ \texttt{TRANSCRIBIENDO} $\rightarrow$ \texttt{GUARDANDO}
$\rightarrow$ \texttt{EMBEDDINGS} $\rightarrow$ \texttt{ANALIZANDO} $\rightarrow$ \texttt{COMPLETADO}
\end{center}

El trabajo puede terminar también como \texttt{CANCELADO} o \texttt{ERROR}. Si el usuario cancela,
el sistema interrumpe el hilo y elimina los datos parciales ya guardados, dejando la base de datos
en un estado consistente.

\section{Módulos funcionales}

\subsection{Transcripción}

El sistema intenta primero obtener los subtítulos del vídeo con \texttt{yt-dlp} (prefiriendo
español y luego inglés). Si existen y tienen calidad suficiente, se usan directamente y se omite
Whisper. En caso contrario, se descarga el audio y se transcribe con Whisper. Cada vía queda
registrada en el campo \texttt{fuenteTranscripcion} (\texttt{YOUTUBE} o \texttt{WHISPER}).

\subsection{Indexación y búsqueda semántica}

La transcripción se divide en fragmentos con sus tiempos. Para cada fragmento se genera un
\textit{embedding} que se almacena con \texttt{pgvector}. Ante una consulta, se calcula el
\textit{embedding} de la pregunta y se recuperan los fragmentos más cercanos por distancia del
coseno.

\subsection{Análisis automático: capítulos y conceptos}

Tras generar los \textit{embeddings}, un servicio de análisis solicita al modelo de lenguaje la
segmentación de la clase en capítulos y la extracción de conceptos clave. Para evitar que el
modelo invente marcas de tiempo, los capítulos se anclan a fragmentos reales de la transcripción
mediante coincidencia léxica de una frase clave, y después se normalizan temporalmente (orden,
eliminación de solapamientos y ajuste de los tiempos de fin). Si el modelo falla, se aplica una
segmentación local por bloques temporales para no interrumpir el procesamiento.

\subsection{Asistente conversacional (RAG)}

El asistente responde preguntas sobre el vídeo combinando recuperación semántica y generación de
lenguaje. Distingue entre preguntas globales (resumen, ideas principales), para las que construye
contexto con el título, los capítulos, los conceptos y una muestra representativa de fragmentos,
y preguntas concretas, para las que recupera los fragmentos más relevantes y cita el momento
exacto del vídeo.

El chat mantiene una \textbf{memoria corta} del diálogo en el \textit{frontend} (las últimas
interacciones, sin persistencia en el servidor). Esa memoria, junto con la última entidad
mencionada, se envía al \textit{backend} para resolver referencias implícitas: por ejemplo, si el
alumno pregunta ``¿y el móvil?'' tras hablar del ``iPhone 17 Pro Max'', el sistema enriquece la
búsqueda con esa entidad para recuperar los fragmentos correctos.

\subsection{Organización y clasificación automática de asignaturas}

Las clases se agrupan en asignaturas. Para reducir el trabajo manual, al terminar el análisis el
sistema sugiere automáticamente una asignatura siguiendo este orden de decisión:

\begin{enumerate}
  \item \textbf{Por canal de YouTube:} si ya existe una asignatura asociada al mismo canal
        (por identificador o por nombre normalizado), se asigna esa.
  \item \textbf{Por similitud semántica:} en caso contrario, se compara el \textit{embedding} del
        vídeo (título, resumen, conceptos y capítulos) con el de cada asignatura existente; si la
        similitud supera un umbral, se reutiliza esa asignatura.
  \item \textbf{Nueva asignatura:} si ninguna resulta adecuada, se crea una nueva (por ejemplo,
        ``Vídeos de \textit{<canal>}'').
\end{enumerate}

En todos los casos, la asignatura se marca como \textbf{sugerencia} (metadato visual): la interfaz
muestra ``Asignatura: X (sugerencia)''. La palabra ``sugerencia'' nunca forma parte del nombre
real. Si el usuario cambia la asignatura a mano, deja de ser sugerencia y el sistema registra que
la asignación fue manual; además, la asignatura ``aprende'' el canal del vídeo para clasificar
mejor los siguientes.

\section{Diseño de la API REST}

La comunicación entre \textit{frontend} y \textit{backend} se realiza mediante una API REST. La
tabla~\ref{tab:endpoints} recoge los \textit{endpoints} más representativos.

\begin{longtable}{p{1.5cm}p{6.5cm}p{5cm}}
  \caption{\textit{Endpoints} principales de la API.}\label{tab:endpoints}\\
  \toprule
  \textbf{Método} & \textbf{Ruta} & \textbf{Descripción} \\
  \midrule
  \endhead
  POST & \texttt{/api/transcripciones/youtube} & Iniciar análisis (asíncrono). \\
  GET & \texttt{/api/transcripciones/youtube/\{id\}} & Estado del trabajo. \\
  GET & \texttt{/api/transcripciones} & Historial de clases. \\
  GET & \texttt{/api/transcripciones/\{id\}} & Detalle de una clase. \\
  PATCH & \texttt{/api/transcripciones/\{id\}/metadata} & Editar metadatos (asignatura, profesor\ldots). \\
  GET & \texttt{/api/transcripciones/\{id\}/capitulos} & Capítulos de la clase. \\
  GET & \texttt{/api/transcripciones/\{id\}/conceptos} & Conceptos clave. \\
  GET & \texttt{/api/transcripciones/\{id\}/buscar} & Búsqueda semántica. \\
  POST & \texttt{/api/transcripciones/\{id\}/conversar} & Chat conversacional con memoria corta. \\
  GET & \texttt{/api/asignaturas} & Listar asignaturas. \\
  \bottomrule
\end{longtable}

\chapter{Pruebas y resultados}
\label{cap:pruebas}

Este capítulo describe las pruebas realizadas para validar el sistema y los resultados
obtenidos, relacionándolos con los objetivos planteados en la introducción.

\section{Pruebas automáticas}

El \textit{backend} cuenta con una batería de pruebas automáticas (JUnit y \texttt{@QuarkusTest})
que se ejecutan sin depender de servicios externos. Para ello, en el entorno de pruebas se usa una
base de datos H2 en memoria y se desactivan las llamadas reales a OpenAI y a \texttt{yt-dlp}
mediante propiedades de configuración. Esto permite verificar la lógica del sistema de forma
rápida y reproducible.

Las pruebas cubren, entre otros aspectos, la validación de URLs de YouTube, el comportamiento del
procesamiento simulado, la cancelación de trabajos y las consultas de capítulos y conceptos. En el
estado actual, la suite se compone de 19 pruebas que se ejecutan correctamente.

% TODO: añadir tabla con el listado de pruebas (clase, qué valida, resultado).

\section{Validación funcional}

Además de las pruebas automáticas, se ha validado el flujo completo de forma manual: procesar un
vídeo de YouTube, comprobar la generación de capítulos y conceptos con marcas de tiempo
coherentes, realizar preguntas al asistente (globales y concretas) y verificar los saltos al
minuto exacto del vídeo. También se han validado los escenarios de la clasificación automática de
asignaturas:

\begin{itemize}
  \item Vídeo de un canal ya asociado a una asignatura $\rightarrow$ se asigna por canal.
  \item Vídeo de un canal nuevo con temática parecida a una asignatura existente $\rightarrow$ se
        asigna por similitud semántica.
  \item Vídeo sin coincidencias $\rightarrow$ se crea una asignatura nueva.
  \item Cambio manual de la asignatura $\rightarrow$ desaparece la marca de ``sugerencia'' y
        persiste tras recargar.
\end{itemize}

\section{Medición de tiempos}

Para localizar los cuellos de botella del procesamiento, la clasificación automática registra el
tiempo de cada fase (obtención de metadatos del canal, coincidencia por canal, generación del
\textit{embedding} del vídeo, carga o generación de \textit{embeddings} de asignaturas, cálculo de
similitud y asociación final). Estas medidas permiten decidir dónde optimizar antes de modificar la
lógica.

% TODO: incluir una tabla/gráfico con tiempos reales medidos sobre varios vídeos (canal, semántica
% y nueva) una vez ejecutado en el entorno con base de datos y clave de OpenAI.

\section{Discusión de resultados}

Los resultados confirman que el sistema cumple los objetivos planteados: transcribe el vídeo,
lo indexa, genera capítulos y conceptos, responde preguntas con referencias al minuto exacto y
sugiere automáticamente la asignatura. Las principales fuentes de latencia son las llamadas a
servicios externos (descarga con \texttt{yt-dlp} y peticiones a OpenAI), lo que orienta el trabajo
futuro hacia su optimización (véase el capítulo~\ref{cap:conclusiones}).

\chapter{Conclusiones y líneas futuras}
\label{cap:conclusiones}

\section{Conclusiones}

El proyecto ha alcanzado su objetivo general: se ha diseñado e implementado una aplicación web
que transforma vídeos de YouTube en material de estudio estructurado y consultable. Repasando los
objetivos específicos planteados en la introducción:

\begin{itemize}
  \item Se obtiene la transcripción de forma eficiente, priorizando los subtítulos de YouTube y
        recurriendo a Whisper solo cuando es necesario.
  \item El contenido se indexa mediante \textit{embeddings}, lo que habilita la búsqueda semántica
        por fragmentos.
  \item Se generan automáticamente capítulos y conceptos clave con marcas de tiempo coherentes.
  \item El asistente conversacional responde preguntas sobre el vídeo con memoria corta del
        diálogo y referencias al minuto exacto.
  \item Las clases se organizan en asignaturas y el sistema sugiere automáticamente la asignatura
        de cada vídeo nuevo.
  \item La interfaz integra reproductor, capítulos, conceptos y chat en una sola pantalla.
\end{itemize}

Más allá del cumplimiento de los requisitos, el trabajo me ha permitido integrar en un sistema
real un conjunto de tecnologías actuales —transcripción automática, \textit{embeddings},
búsqueda semántica y modelos de lenguaje— y enfrentarme a problemas propios de la ingeniería del
software, como el procesamiento asíncrono, la degradación controlada ante fallos de servicios
externos y la consistencia de los datos.

\subsection{Orientación social}

ComprendiA contribuye a un acceso más equitativo al conocimiento: permite navegar las clases por
conceptos en lugar de verlas completas, favorece el aprendizaje autónomo y ayuda a estudiantes
con ritmos o necesidades distintas. En este sentido, el proyecto se alinea con los Objetivos de
Desarrollo Sostenible relativos a la educación de calidad (ODS~4), la reducción de las
desigualdades (ODS~10) y la innovación (ODS~9).

\section{Líneas futuras}

El sistema es un prototipo funcional con margen claro de evolución. Las principales líneas de
trabajo futuro son:

\begin{itemize}
  \item \textbf{Integración con plataformas educativas (LMS).} Conectar ComprendiA con Moodle,
        Blackboard u otros sistemas para procesar directamente las clases del entorno
        universitario.
  \item \textbf{Asistente en tiempo real.} Llevar también el chat a WebSockets, con respuesta en
        \textit{streaming} (token a token), para una experiencia conversacional más inmediata. El
        progreso del análisis ya se transmite por WebSocket; faltaría extenderlo al asistente.
  \item \textbf{Gestión multiusuario.} Añadir autenticación y perfiles para soportar varios
        estudiantes y profesores.
  \item \textbf{Optimización del procesamiento.} Reducir la latencia de las llamadas externas
        (reutilizar la consulta de \texttt{yt-dlp}, cachear \textit{embeddings} de asignaturas,
        paralelizar fases) a partir de los tiempos medidos.
  \item \textbf{Mejora de la clasificación.} Umbral de similitud configurable, sugerencia también
        del profesor y refinamiento del aprendizaje del canal.
  \item \textbf{Reproductor enriquecido.} Marcar sobre la barra de progreso los capítulos y
        fragmentos analizados, sincronizando la posición del vídeo con la transcripción.
\end{itemize}

% TODO: si procede, mencionar el contexto del programa TALENT 25-26 y el recorrido del proyecto.


% ====================  BIBLIOGRAFÍA Y ANEXOS  ===============================
\bibliographystyle{plain}
\bibliography{bibliografia}
\addcontentsline{toc}{chapter}{Bibliografía}

\appendix
\chapter{Manual de instalación y uso}
\label{anexo:manual}

Este anexo resume cómo instalar, arrancar y utilizar ComprendiA. La información detallada y
actualizada se encuentra en el fichero \texttt{README.md} del repositorio del proyecto.

\section{Requisitos previos}

\begin{itemize}
  \item \texttt{yt-dlp} y \texttt{ffmpeg} instalados en el sistema.
  \item Java~21 y Maven.
  \item Node.js y Angular CLI.
  \item Una base de datos PostgreSQL (por ejemplo, en Neon) y una clave de la API de OpenAI.
\end{itemize}

\section{Configuración}

Las credenciales se definen en el fichero \texttt{backend/.env} (no versionado), con la clave de
OpenAI y la cadena de conexión a la base de datos. Existe una plantilla \texttt{.env.example} como
referencia.

\section{Arranque}

\begin{enumerate}
  \item \textbf{Backend:} desde \texttt{backend/}, ejecutar \texttt{mvn quarkus:dev}. El servidor
        escucha en \texttt{http://localhost:8080}.
  \item \textbf{Frontend:} desde \texttt{frontend/}, ejecutar \texttt{npm install} y \texttt{ng
        serve}. La aplicación se abre en \texttt{http://localhost:4200}.
\end{enumerate}

\section{Uso básico}

\begin{enumerate}
  \item Pegar la URL de un vídeo de YouTube y pulsar \textit{Analizar}.
  \item Seguir el progreso del procesamiento (descarga, transcripción, guardado,
        \textit{embeddings} y análisis).
  \item Al terminar, se abre la pantalla de la clase con el reproductor, los capítulos, los
        conceptos clave y el asistente.
  \item Editar la asignatura, el profesor o la fecha; usar el chat para preguntas globales o
        concretas; y saltar al minuto exacto desde capítulos, conceptos o respuestas.
\end{enumerate}

% TODO: añadir capturas de pantalla de la pantalla de clase, del chat y del grid de asignaturas.


\end{document}
